% 03-methods — v1.0
\subsection*{Datasets}
\textbf{EEG Motor Movement/Imagery Dataset} \citep{schalk2009eegbci}. All 109 subjects were used: 64 channels at 160\,Hz, runs 1--2 for the
eyes-open and eyes-closed baselines and runs 4, 8 and 12 for imagined left/right hand
movement, as defined in the dataset documentation. Motor decoding uses the ten subjects for whom the imagery runs
were degraded in this pass, and is treated as exploratory.

\textbf{Affective dataset} \citep{sanchezreolid2022emotion,sanchezreolid2025dataset}.
The open CC BY 4.0 dataset from Zenodo 14787743 contains EEG from 20 participants who
viewed International Affective Picture System images selected to induce four
valence--arousal quadrants, recorded with a low-cost Emotiv EPOC+ and a high-end
BrainVision system. The archive's \texttt{EEG/Experiment\_files} directory contains 18 BrainVision
recordings (66 channels, 500\,Hz, sampling interval 2000\,\textmu s); a nineteenth
recording sits in a separate \texttt{EEG registro\_nuevo} folder and was not used. All
18 recordings in the analysed directory completed the pipeline. The files are named
\texttt{eegbci<subject>-<session>} by the archive; this is the archive's own naming for
the affective recordings and does not refer to the EEGBCI motor dataset. The
valence/arousal labels come from the affective picture protocol, not from a motor task. BrainVision
marker codes map to the quadrants as follows: 1 = high valence/high arousal,
2 = low valence/high arousal, 3 = high valence/low arousal, 4 = low valence/low
arousal, 10/14 = neutral. Each recording contains approximately 10 trials per quadrant
and 8 neutral trials; neutral trials are excluded from valence and arousal analyses.
Quadrant labels are the normative assignments of the International Affective Picture
System \citep{lang2008iaps}, not participants' own ratings.
Marker positions are used as onsets and each trial is epoched 0--2.5\,s after onset.
The low-cost Emotiv arm is analysed for the 12 participants whose CSV files are
usable in the distributed archive (five files are empty or header-only and three fail
the timestamp alignment). The device has been used in validation work on
event-related potentials \citep{badcock2015emotiv}. The files contain 14 EEG channels at
128\,Hz in device units with absolute wall-clock timestamps. The log session whose
absolute trial times fall inside the CSV span is selected; trials are aligned by
converting each log onset to an absolute time (SessionStartDateTimeUtc plus the
millisecond onset) and matching it to the nearest CSV sample, with a clock offset chosen
from $\{0,\pm1,\pm2\}$ hours by maximising the number of trials inside the recording
span. Emotiv epochs run from $-0.2$ to $+2.3$\,s around the aligned sample (the
BrainVision epochs run 0--2.5\,s), and the same band-power and LDA pipeline is applied.
The median absolute difference between a log onset and its nearest CSV sample was
25\,ms (95th percentile 59\,ms, maximum 70\,ms; stored per subject in
\texttt{results/emotiv.json}), and no trial exceeded the 1\,s alignment screen (this is a timestamp screen, not an amplitude rejection threshold). The
two condition arms in the archive are not otherwise distinguished. This yields a
three-way comparison: high-end full montage, low-cost 14-channel headset, and high-end
data reduced to a single degraded channel.

Both datasets are de-identified and openly available. No new data were collected.

\subsection*{Degradation grid}
EEGBCI recordings are processed at a full-montage reference and at twelve
single-channel conditions: a cumulative acquisition chain (Fz; +128\,Hz; +1\,\textmu V
quantisation; +1--40\,Hz band-limiting), a noise arm (20/10/5/0\,dB pink noise, no
blinks; pink noise is a conservative broadband proxy and real consumer noise also
contains white electronics noise, mains interference and drift, so the noise arm should
be read as one plausible contamination profile rather than a device model), a blink arm (5/15/30 per minute, no noise) and a combined arm (0\,dB + 30
blinks/min). The single electrode is selected by canonical name matching because
EEGBCI channel labels carry punctuation (e.g.\ \texttt{Fz.}); the chosen label is
written to the condition metadata. Blink transients are 350\,ms Hann-windowed 2\,Hz
templates scaled to four standard deviations of the channel; with one channel the
template cannot be given a frontal topography, which is a limitation of the model.

For the affective dataset the same single-channel chain (128\,Hz, 1\,\textmu V,
1--40\,Hz, Fz) is applied to the high-end recording and compared with the full
montage. The noise and blink arms are not applied to the affective data; that
comparison is left to future work.

\subsection*{Site-matched sensitivity analysis}
A single Fz electrode and a 64-channel average are not the same measurement: alpha is
posterior-dominant, so comparing them confounds channel count with electrode site. We
therefore compute, in the eyes-closed run, Fz power from the full-montage recording,
Fz power after the full single-channel chain, the 64-channel average and the average
over posterior channels. The site-matched ratio (single-channel Fz versus full-montage
Fz) is the primary estimate of what the acquisition chain costs; the scalp-average
comparison is reported as the naive estimate that a reader would get without
site-matching.

\subsection*{Analyses}
\textbf{A. Alpha reactivity.} Band power (8--12\,Hz) is computed with Welch's method
(2\,s windows, 50\% overlap) on continuous data; the outcome is the natural-log ratio
of eyes-closed to eyes-open power, the classical resting-state contrast
\citep{barry2007eyes}. The primary retention estimator is the geometric mean ratio derived from the paired
subject-level log difference. For the acquisition chain the reference is the same-site
full recording (Fz in the full montage); for the noise, blink and combined arms the
reference is the clean single-channel chain, because those manipulations are
incremental contamination of that chain. A degradation claim is reported as
\emph{supported} when the \emph{lower} 95\% confidence bound of the relevant retention
ratio exceeds 0.50; as \emph{not supported} when the lower bound is at or below 0.50;
and the decision matrix reports the retention interval alongside the verdict. The 50\% threshold is a planned
convention of this benchmark, not a literature-derived constant, and it is not a formal
non-inferiority test. The pre-specified 60\% practical floor for balanced accuracy is
reported separately as a stricter auxiliary criterion. Absolute powers are released so readers can evaluate alternative summaries.

\textbf{B. Motor decoding (exploratory).} Single-trial log band power (8--30\,Hz;
mu/beta event-related desynchronisation \citep{pfurtscheller1999erd}) with shrinkage
LDA and 5-fold stratified within-subject cross-validation. A CSP-based reference
decoder (CSP fitted inside each fold) is reported as a sensitivity analysis. Because
CSP requires multiple channels, the single-channel condition cannot receive the same
spatial filter; that asymmetry is part of the constraint quantified here. The arm is
reported as invalid unless a reference decoder's 95\% CI excludes chance (0.50).
Because nominal chance levels can be exceeded by chance, interpretation uses empirical
confidence intervals rather than the theoretical level
\citep{combrisson2015chance,billinger2012significance}.

\textbf{C. Affective decoding.} For each recording, log band power in theta (4--8\,Hz),
alpha (8--12\,Hz) and beta (13--30\,Hz) is extracted per channel and trial; the
classifier is shrinkage LDA with 5-fold stratified within-subject cross-validation.
Three targets are decoded: valence (quadrants 1/3 versus 2/4), arousal (1/2 versus
3/4) and the four quadrants. The same pipeline is run (i) on the full montage,
(ii) after each cumulative acquisition step (128\,Hz, quantisation, band-limiting),
(iii) on single sites Fz, Cz, Pz and Oz, and (iv) on the full montage with only
theta+alpha features and with only beta features. The last two conditions separate a
representation effect from a band/EMG effect. Chance is 0.50 for the binary targets and
0.25 for the four-class target; the empirical CI is the inferential criterion
\citep{combrisson2015chance}. The low-cost Emotiv arm is analysed separately and compared
with the high-end arm on the 12 participants present in both, using a TOST procedure
for equivalence.

\textbf{D. Reliability.} Alpha power is computed in consecutive 2\,s windows of the
eyes-closed run; odd and even windows are correlated, and the Spearman--Brown formula is
reported as a descriptive \emph{short-lag stability} measure. Odd and even windows are
about two seconds apart, so this is not a split-half reliability coefficient and is not
directly comparable to the ICC. In the ten-subject subset with three repeated
imagined-hand runs, ICC(2,1) \citep{shrout1979icc} of alpha power across runs is reported
as a block-level estimate. Neither is a test--retest coefficient across sessions, which
is not identifiable in these datasets.

\textbf{E. Blink susceptibility of band ratios.} Pre-specified but not implemented;
reported as untested.

\subsection*{Statistical analysis}
Subject-level outcomes are summarised with means and 95\% confidence intervals
(t-distribution, $df=n-1$). Condition effects are reported both as unpaired means and
as paired contrasts against the reference for the same subjects. Benjamini--Hochberg
FDR \citep{benjamini1995fdr} is applied within the alpha condition family and within the
affective family; both raw and adjusted values are released. Mixed-effects models are
not used here, and the exploratory nature of the study is stated explicitly. Equivalence
between the low-cost and high-end devices is tested with a two one-sided tests (TOST)
procedure using a $\pm 0.10$ balanced-accuracy margin
\citep{lakens2018equivalence}.

\subsection*{Reproducibility}
All degradation parameters, random seeds, dataset file hashes (written by
\texttt{scripts/01\_download.py}), analysis scripts and derived feature tables are
included in the repository. Raw datasets are not redistributed; download scripts
reproduce the exact files used. Figure source data are archived with the results JSON.
Artifact handling is deliberately minimal: no Autoreject or ICA is applied, because
single-channel data cannot support ICA and the benchmark is about the acquisition
conditions rather than about cleaning; the ocular-artifact literature
\citep{gratton1983ocular,croft2000ocular} and the EMG-contamination literature
\citep{whitham2007emg,goncharova2003emg,muthukumaraswamy2013highfreq} are discussed as
limitations. Single-channel test--retest work reports eyes-closed relative-power ICCs of
roughly 0.76--0.85 for a consumer device \citep{rogers2016reliability}; our short-lag
estimate is a within-session stability measure and is not directly comparable.
